\documentclass[12pt,a4paper]{article}

\usepackage[a4paper,margin=18mm]{geometry}
\usepackage[T2A]{fontenc}
\usepackage[utf8]{inputenc}
\usepackage[russian]{babel}
\usepackage{amsmath,amssymb,mathtools}
\usepackage{array,booktabs,longtable}
\usepackage{xcolor}
\usepackage{hyperref}
\usepackage{tikz}

\hypersetup{
  colorlinks=true,
  linkcolor=blue!55!black,
  urlcolor=blue!55!black
}

\setlength{\parindent}{0pt}
\setlength{\parskip}{0.6em}
\renewcommand{\arraystretch}{1.2}

\begin{document}

\begin{titlepage}
  \thispagestyle{empty}
  \begin{center}
    \vspace*{0.18\textheight}
    {\Huge\bfseries Ревью домашней работы\par}
    \vspace{2.2cm}
    {\Large Автор: Лёвин Артём Александрович\par}
    \vspace{0.8cm}
    {\large 9 сентября 2026 года\par}
    \vfill
    \begin{tikzpicture}[x=1cm,y=1cm]
      \draw[blue!55!black, line width=0.8pt]
        (0,0) .. controls (1.8,0.55) and (3.8,-0.55) .. (5.6,0)
              .. controls (7.2,0.45) and (8.8,-0.35) .. (10.2,0.05);
    \end{tikzpicture}
  \end{center}
\end{titlepage}

\section*{Сводная таблица}
\small
\setlength{\tabcolsep}{3pt}
\begin{longtable}{@{}p{0.08\linewidth}p{0.17\linewidth}p{0.27\linewidth}p{0.16\linewidth}p{0.22\linewidth}@{}}
\toprule
\textbf{№ задания} & \textbf{Тема} & \textbf{Суть ошибки} & \textbf{Ваш ответ} & \textbf{Правильный ответ} \\
\midrule
\endfirsthead
\toprule
\textbf{№ задания} & \textbf{Тема} & \textbf{Суть ошибки} & \textbf{Ваш ответ} & \textbf{Правильный ответ} \\
\midrule
\endhead
\hyperref[task:3]{3} & Квадратичная функция по графику & Задание не выполнено: коэффициенты и значение функции не найдены. & не указан & \(f(x)=x^2+x-3\), \(f(-1)=-3\) \\
\addlinespace
\hyperref[task:9]{9} & Парабола и прямая: точки пересечения & Начало решения есть, но обе функции не восстановлены и абсциссы пересечений не найдены. & не указан & \(f(x)=x^2+3x-2\), \(g(x)=2x+2\), \(x=\dfrac{-1\pm\sqrt{17}}{2}\) \\
\addlinespace
\hyperref[task:10]{10} & Гипербола и прямая: пересечения & Гипербола восстановлена верно, но формула прямой и точки пересечения не найдены. & не указан & \(f(x)=\dfrac{2}{x}\), \(g(x)=x+1\); точки \((-2,-1)\), \((1,2)\) \\
\bottomrule
\end{longtable}
\normalsize
\setlength{\tabcolsep}{6pt}

\newpage
\thispagestyle{empty}
\vspace*{\fill}
\begin{center}
  {\Huge\bfseries Разбор ошибок}
\end{center}
\vspace*{\fill}
\newpage

\phantomsection\label{task:3}
\section*{Задание 3}

\subsection*{Текст задания}
На рисунке изображена парабола
\[
  f(x)=x^2+bx+c.
\]
Восстановите формулу и найдите \(f(-1)\).

\subsection*{Ваше решение}
Решение задания отсутствует.

\subsection*{Ваш ответ}
Не указан.

\textcolor{red}{Ошибка:} задание не выполнено, поэтому коэффициенты \(b\), \(c\) и требуемое значение \(f(-1)\) не определены.

\subsection*{Где именно возникает ошибка}
В записанном решении нет математических шагов по заданию 3, поэтому последнего выполненного корректного шага нет. Первый неполный шаг --- само начало восстановления формулы: нужно было считать с графика координаты отмеченных точек и подставить их в заданный вид функции.

По графику читаются точки
\[
  A(0,-3), \qquad B(2,3).
\]
Именно эти данные позволяют найти неизвестные коэффициенты.

\subsection*{Почему это неверно}
Требование задания состоит из двух обязательных частей: восстановить формулу функции и затем вычислить её значение при \(x=-1\). Если коэффициенты \(b\) и \(c\) не найдены, формула остаётся неопределённой, а вычисление \(f(-1)\) невозможно выполнить обоснованно.

Здесь коэффициент при \(x^2\) уже задан и равен \(1\), поэтому неизвестны только два числа: \(b\) и \(c\). Для их определения достаточно двух независимых условий, которые дают две точки графика.

\textcolor{blue!60!black}{Ключевая идея:} если точка \((x_0,y_0)\) лежит на графике функции, то её координаты удовлетворяют равенству \(y_0=f(x_0)\). Поэтому координаты отмеченных точек нужно по очереди подставить в формулу \(f(x)=x^2+bx+c\).

\subsection*{Исправление от последнего правильного шага}
Так как записанного хода решения нет, начинаем с первого необходимого шага --- используем точку \(A(0,-3)\):
\[
  -3=0^2+b\cdot0+c.
\]
Отсюда сразу
\[
  c=-3.
\]

Теперь используем точку \(B(2,3)\):
\[
  3=2^2+2b+c.
\]
Подставляем \(c=-3\):
\[
  3=4+2b-3,
\]
\[
  3=1+2b,
\]
\[
  2b=2,
\]
\[
  b=1.
\]

\subsection*{Полное правильное решение}
Получили
\[
  b=1, \qquad c=-3.
\]
Следовательно,
\[
  f(x)=x^2+x-3.
\]

Теперь вычислим требуемое значение:
\[
  f(-1)=(-1)^2+(-1)-3.
\]
Так как \((-1)^2=1\), имеем
\[
  f(-1)=1-1-3=-3.
\]

\subsection*{Проверка}
Проверим восстановленную формулу по обеим отмеченным точкам.

Для \(A(0,-3)\):
\[
  f(0)=0^2+0-3=-3.
\]
Координата точки совпадает.

Для \(B(2,3)\):
\[
  f(2)=2^2+2-3=4+2-3=3.
\]
Координата точки также совпадает. Значит, коэффициенты восстановлены согласованно с графиком.

\textcolor{teal!60!black}{Итог:}
\[
  \boxed{f(x)=x^2+x-3}, \qquad \boxed{f(-1)=-3}.
\]

\subsection*{Задача на закрепление}
Парабола \(y=x^2+bx+c\) проходит через точки \((0,-2)\) и \((3,4)\). Восстановите формулу функции и найдите её значение при \(x=-2\).

\newpage
\phantomsection\label{task:9}
\section*{Задание 9}

\subsection*{Текст задания}
На рисунке изображены графики
\[
  f(x)=x^2+bx+c, \qquad g(x)=mx+n.
\]
Восстановите обе функции и найдите все абсциссы точек пересечения.

\subsection*{Ваше решение}
Записаны общие виды функций
\[
  f(x)=x^2+bx+c, \qquad g(x)=mx+n,
\]
а для параболы отмечены точки
\[
  A(0,-2), \qquad B(1,2).
\]
Дальнейшие вычисления отсутствуют.

\subsection*{Ваш ответ}
Не указан.

\textcolor{red}{Ошибка:} решение остановлено до определения коэффициентов функций и до решения уравнения их пересечения.

\subsection*{Где именно возникает ошибка}
Последний корректный шаг --- распознаны общий вид параболы и прямой, а также две точки параболы \(A(0,-2)\) и \(B(1,2)\).

Первый неполный шаг возникает сразу после этого: координаты точек не подставлены в формулу параболы, поэтому \(b\) и \(c\) не найдены. Формула прямой также не восстановлена, и равенство \(f(x)=g(x)\) не составлено.

\subsection*{Почему это неверно}
Само указание двух точек ещё не завершает восстановление функции. Коэффициенты должны быть получены из условий принадлежности точек графику.

Для параболы неизвестны \(b\) и \(c\), поэтому нужны два уравнения. Для прямой неизвестны \(m\) и \(n\), поэтому также нужны два независимых условия. На графике для прямой читаются точки
\[
  C(-1,0), \qquad D(1,4).
\]
После восстановления обеих формул необходимо найти именно точки пересечения графиков. Их абсциссы удовлетворяют уравнению
\[
  f(x)=g(x).
\]
Без этого шага последняя часть задания остаётся невыполненной.

\textcolor{blue!60!black}{Ключевая идея:} сначала восстановить каждую функцию независимо по её отмеченным точкам, а затем приравнять полученные выражения. Равенство ординат в точке пересечения превращает графическую задачу в обычное квадратное уравнение.

\subsection*{Исправление от последнего правильного шага}
Начинаем с уже верно выписанных точек параболы.

Из точки \(A(0,-2)\):
\[
  -2=0^2+b\cdot0+c,
\]
поэтому
\[
  c=-2.
\]

Из точки \(B(1,2)\):
\[
  2=1^2+b\cdot1+c.
\]
Подставляем \(c=-2\):
\[
  2=1+b-2,
\]
\[
  2=b-1,
\]
\[
  b=3.
\]
Следовательно,
\[
  f(x)=x^2+3x-2.
\]

Теперь восстановим прямую по точкам \(C(-1,0)\) и \(D(1,4)\). Её угловой коэффициент:
\[
  m=\frac{4-0}{1-(-1)}=\frac{4}{2}=2.
\]
Используем точку \(C(-1,0)\):
\[
  0=2\cdot(-1)+n,
\]
откуда
\[
  n=2.
\]
Значит,
\[
  g(x)=2x+2.
\]

\subsection*{Полное правильное решение}
Точки пересечения имеют одинаковую ординату, поэтому решаем уравнение
\[
  x^2+3x-2=2x+2.
\]
Перенесём всё в левую часть:
\[
  x^2+x-4=0.
\]

Вычислим дискриминант:
\[
  D=1^2-4\cdot1\cdot(-4)=1+16=17.
\]
Тогда
\[
  x_{1,2}=\frac{-1\pm\sqrt{17}}{2}.
\]
То есть
\[
  x_1=\frac{-1-\sqrt{17}}{2},
  \qquad
  x_2=\frac{-1+\sqrt{17}}{2}.
\]

\subsection*{Проверка}
Разность функций равна
\[
  f(x)-g(x)=x^2+x-4.
\]
Оба найденных значения \(x\) являются корнями уравнения \(x^2+x-4=0\), следовательно, для каждого из них выполняется \(f(x)=g(x)\).

Для дополнительной проверки оценим корни приближённо:
\[
  x_1\approx -2{,}56, \qquad x_2\approx 1{,}56.
\]
По рисунку одна точка пересечения действительно расположена левее оси \(Oy\), а вторая --- правее, что согласуется с найденными значениями.

\textcolor{teal!60!black}{Итог:}
\[
  \boxed{f(x)=x^2+3x-2}, \qquad
  \boxed{g(x)=2x+2},
\]
\[
  \boxed{x=\frac{-1-\sqrt{17}}{2}}, \qquad
  \boxed{x=\frac{-1+\sqrt{17}}{2}}.
\]

\subsection*{Задача на закрепление}
Парабола \(f(x)=x^2+bx+c\) проходит через точки \((0,-1)\) и \((2,5)\), а прямая \(g(x)=mx+n\) --- через точки \((-1,0)\) и \((1,2)\). Восстановите обе функции и найдите все абсциссы точек их пересечения.

\newpage
\phantomsection\label{task:10}
\section*{Задание 10}

\subsection*{Текст задания}
На рисунке изображены графики
\[
  f(x)=\frac{k}{x}+a, \qquad g(x)=mx+n.
\]
Требуется:
\begin{enumerate}
  \item восстановить обе формулы;
  \item найти абсциссы точек пересечения;
  \item найти координаты точек пересечения.
\end{enumerate}

\subsection*{Ваше решение}
Записано:
\[
  a=0, \qquad f(x)=\frac{k}{x}+a,
\]
\[
  B(2,1) \Rightarrow k=2,
\]
поэтому получено
\[
  f(x)=\frac{2}{x}.
\]
Для прямой записан общий вид
\[
  g(x)=mx+n
\]
и отмечена точка \(H(-1,0)\). На этом решение остановлено.

\subsection*{Ваш ответ}
Не указан.

\textcolor{red}{Ошибка:} первая функция восстановлена, но формула прямой и все требуемые данные о точках пересечения не найдены.

\subsection*{Где именно возникает ошибка}
Последний корректный шаг --- получена формула
\[
  f(x)=\frac{2}{x},
\]
а также распознана точка прямой \(H(-1,0)\).

Первый неполный шаг --- попытка восстановить \(g(x)=mx+n\) остановлена после одной точки. Одной точки недостаточно для однозначного определения двух коэффициентов \(m\) и \(n\). На графике есть ещё одна удобная точка прямой:
\[
  D(2,3).
\]
После определения прямой необходимо дополнительно решить уравнение пересечения и найти ординаты, чего в записи нет.

\subsection*{Почему это неверно}
Прямая \(g(x)=mx+n\) определяется двумя независимыми условиями. Одна точка даёт только одно уравнение с двумя неизвестными. Поэтому запись \(H(-1,0)\) сама по себе не позволяет получить ни \(m\), ни \(n\) однозначно.

Кроме того, по условию недостаточно восстановить функции: нужно найти и абсциссы, и полные координаты всех точек пересечения. Для этого после нахождения \(g(x)\) следует решить
\[
  \frac{2}{x}=g(x),
\]
обязательно учитывая ограничение \(x\ne0\).

\textcolor{blue!60!black}{Ключевая идея:} использовать две точки прямой для нахождения \(m\) и \(n\), затем приравнять значения функций. При работе с дробью \(\frac{2}{x}\) перед умножением уравнения на \(x\) отдельно фиксируется условие \(x\ne0\).

\subsection*{Исправление от последнего правильного шага}
Формула гиперболы уже получена верно:
\[
  f(x)=\frac{2}{x}.
\]

Восстановим прямую по точкам \(H(-1,0)\) и \(D(2,3)\). Угловой коэффициент равен
\[
  m=\frac{3-0}{2-(-1)}=\frac{3}{3}=1.
\]
Подставим точку \(H(-1,0)\) в \(g(x)=x+n\):
\[
  0=-1+n,
\]
поэтому
\[
  n=1.
\]
Следовательно,
\[
  g(x)=x+1.
\]

\subsection*{Полное правильное решение}
Ищем точки пересечения:
\[
  \frac{2}{x}=x+1,
  \qquad x\ne0.
\]
Умножаем обе части на \(x\). Это допустимо, поскольку \(x=0\) заранее исключён из области определения:
\[
  2=x^2+x.
\]
Перенесём всё в одну часть:
\[
  x^2+x-2=0.
\]
Разложим на множители:
\[
  (x+2)(x-1)=0.
\]
Отсюда
\[
  x_1=-2, \qquad x_2=1.
\]
Оба значения отличны от нуля, поэтому оба допустимы.

Теперь найдём ординаты.

При \(x=-2\):
\[
  y=f(-2)=\frac{2}{-2}=-1.
\]
Получаем точку
\[
  (-2,-1).
\]

При \(x=1\):
\[
  y=f(1)=\frac{2}{1}=2.
\]
Получаем точку
\[
  (1,2).
\]

\subsection*{Проверка}
Проверим обе точки по прямой \(g(x)=x+1\).

Для \((-2,-1)\):
\[
  g(-2)=-2+1=-1.
\]
Значение совпадает с \(f(-2)\).

Для \((1,2)\):
\[
  g(1)=1+1=2.
\]
Значение совпадает с \(f(1)\).

Значит, обе найденные точки действительно принадлежат и гиперболе, и прямой.

\textcolor{teal!60!black}{Итог:}
\[
  \boxed{f(x)=\frac{2}{x}}, \qquad \boxed{g(x)=x+1},
\]
\[
  \boxed{x=-2,\;1},
\]
\[
  \boxed{(-2,-1),\;(1,2)}.
\]

\subsection*{Задача на закрепление}
График функции \(f(x)=\dfrac{k}{x}+a\) имеет горизонтальную асимптоту \(y=1\) и проходит через точку \((2,3)\). Прямая \(g(x)=mx+n\) проходит через точки \((-1,0)\) и \((2,3)\). Восстановите обе функции и найдите координаты всех точек их пересечения.

\end{document}
